\scanpage{6}
$H^i(k,\unsure{G|C})$ provient d'un $H^i(k,{}_{\unsure{n|-}}G)$. Or si $G$ est une VA ou un tore, ${}_nG$ est fini sur $k$. Cela montre que les conditions du lemme 1 sont vérifiées pour a) b) c). Reste à vérifier d).

Pour \unsure{ceci|cela}, on va \deleted{\unsure{supposer}} utiliser \deleted{\unsure{un}} le fait que $H^2(k,U)$ devient nul sur une extension de $k$ de la forme $k^{p^{-n}}$, [car $U$ \unsure{splitte} sur une telle extension en groupes additifs] et on va prouver le

\uline{Lemme 2} Pour tt schéma en groupes $G$ sur $k$, le noyau de $H^i(k,G)-H^i(k^{p^{-n}},G)$ est égal à celui de
\[
H^i(k,G)\xrightarrow{F^n}H^i(k,G^{(p^n)})
\]
défini par $F^n$ ($F=\mathrm{Frob}$). Plus précisément, on a commutativité dans le diagramme
% slanted
\margin{\unsure{Pourrait} s'utiliser directement, $+$ \unsure{cette} \unsure{déduction} \unsure{de} \unsure{truc} \unsure{pourrait} \unsure{ramener} au cas du groupe $K$. \deleted{\unsure{faudrait}} vérifier ???}

\deleted{[N.B.}
\diagram{\begin{tikzcd}[ampersand replacement=\&]
H^i(k,G) \arrow[r] \arrow[dr,"F^n"'] \& H^i(k^{p^{-n}},G)\deleted{\simeq H^i(k,G^{(p^n)})} \arrow[d,dash,"\wr"] \\
\& H^i(k,G^{(p^n)})
\end{tikzcd}}
\diagram{\begin{tikzcd}[ampersand replacement=\&,column sep=small]
G \arrow[d,"p"'] \& G^{(p^n)} \arrow[l,"p^n"'] \arrow[d] \arrow[dl,phantom,"\text{cart.}" description] \& G \arrow[l,"F^n"'] \arrow[ll,bend right=38,"\unsure{\varphi_G|\phi_G}"'] \arrow[dl,"p"] \\
\begin{gathered}S\\ \Vert\\ \operatorname{Sp}(k)\end{gathered} \& \begin{gathered}S\\ \Vert\\ \operatorname{Sp}(k)\end{gathered} \arrow[l,"\varphi_S"] \&
\end{tikzcd}}
\[
\lambda\longmapsto\lambda^{p^n}
\]
Pour le voir, on est ramené au cas du $H^0$, où c'est trivial [N.B. On n'a pas besoin que $G$ soit \unsure{représentable}\ldots].
% archive sheet number: 5

